\scanpage{179}\mspage{34}
qui sont des variétés à bord, les \deleted{\unsure{morp|morph}} flèches du diagramme étant pour l'essentiel des fibrations loc. triviales de variétés à bord sur d'autres (à fibres des variétés à bord), \uline{et} \deleted{\unsure{les|des}}\unsure{des|les} inclusions des bords dans les variétés à bord \deleted{(}[en fait, on trouve des inclusions légèrement plus générales, certains des bords qui apparaissent étant munis d'une décomposition cellulaire ``élémentaire'', \deleted{\unsure{en cellules|en éléments}}\inserted{i.e. \unsure{dont}} \unsure{toutes} les strates \inserted{fermées} sont encore des variétés à bord qui se recollent suivant des parties communes du bord~; et il y a lieu d'envisager \inserted{aussi} les inclusions de ces morceaux les uns dans les autres\ldots.\deleted{)}], et se réconstitue à partir de ce diagramme \inserted{par recollement}. En d'autres termes, on a une description-dévissage canonique des espaces modérés compacts \inserted{$X$}, éventuellement munis de stratifications équisingulières à strates non sing., en termes de diagrammes \inserted{finis} de nature précisée formés avec des variétés à bord. Quand on s'intéresse aux faisceaux \inserted{(ou complexes de faisceaux, ou $n$-champs\ldots)} $\unsure{P|\mathcal{P}}$-constructibles sur $X$, où $P$ est une telle stratification fixée, ceux-ci s'explicitent en termes du diagramme envisagé, dont il
% archive sheet number: 179
